\documentclass[11pt,a4paper]{article}
\usepackage[utf8]{inputenc}
\usepackage[T1]{fontenc}
\usepackage[catalan]{babel}
\usepackage{lmodern}
\usepackage[sfdefault,lining]{FiraSans}
\usepackage[a4paper,top=2.3cm,bottom=2.2cm,left=1.8cm,right=1.8cm,headheight=1.3cm,headsep=0.4cm,footskip=1cm]{geometry}
\usepackage{amsmath,amssymb}
\usepackage{icomma}
\usepackage{graphicx}
\usepackage[table]{xcolor}
\usepackage{tikz}
\usetikzlibrary{calc}
\usepackage{fontawesome5}
\usepackage{ccicons}
\usepackage{fancyhdr}
\usepackage{lastpage}
\usepackage{tasks}
\usepackage{needspace}
\usepackage[most]{tcolorbox}
\usepackage{array}
\usepackage[hidelinks,pdftitle={Fitxa d'exercicis de la unitat 3: Sistemes d'equacions i inequacions. 1r de Batxillerat MI},pdfsubject={Matemàtiques I},pdfkeywords={sistemes d'equacions, Gauss, sistemes no lineals, inequacions, sistemes d'inequacions, regions del pla}]{hyperref}

\definecolor{navy}{HTML}{1B2748}
\definecolor{accent}{HTML}{E8603A}
\definecolor{teal}{HTML}{1F9E90}
\definecolor{soft}{HTML}{EEF1F8}
\definecolor{muted}{HTML}{5D6685}

\setlength{\parindent}{0pt}
\setlength{\parskip}{3pt}
\settasks{label=\alph*), label-format=\bfseries\color{navy}, label-width=1.3em, item-indent=1.8em, column-sep=0.6em, after-item-skip=5pt, before-skip=3pt, after-skip=4pt}

% ---------- capçalera i peu ----------
\newcommand{\llicencia}{\href{https://creativecommons.org/licenses/by-nc-sa/4.0/deed.ca}{\ccbyncsa}\ \ Llicència Creative Commons Reconeixement-NoComercial-CompartirIgual 4.0 Internacional (CC BY-NC-SA 4.0)}
\pagestyle{fancy}
\fancyhf{}
\fancyhead[L]{\raisebox{-0.35\height}{\includegraphics[height=0.95cm]{logo-gva.png}}}
\fancyhead[C]{\small\color{muted}Fitxa d'exercicis · Unitat 3: Sistemes d'equacions i inequacions · 1r Batxillerat MI}
\fancyhead[R]{\raisebox{-0.35\height}{\includegraphics[height=0.95cm]{logo-vermellar.png}}}
\fancyfoot[L]{\scriptsize\color{muted}\llicencia}
\fancyfoot[R]{\small\color{muted}\thepage}
\renewcommand{\headrulewidth}{0.4pt}
\renewcommand{\headrule}{\color{muted!40}\hrule height \headrulewidth}
\fancypagestyle{primera}{\fancyhf{}\renewcommand{\headrulewidth}{0pt}
  \fancyfoot[L]{\scriptsize\color{muted}\llicencia}
  \fancyfoot[R]{\small\color{muted}\thepage}}

% ---------- nivells i exercicis ----------
\newcommand{\estrelles}[1]{{\color{accent}\ifcase#1\or\faStar\faStar[regular]\faStar[regular]\or\faStar\faStar\faStar[regular]\or\faStar\faStar\faStar\fi}}
\newcounter{exer}
\newcommand{\ex}[1]{\par\needspace{4\baselineskip}\medskip\refstepcounter{exer}\noindent{\color{navy}\bfseries\theexer.}\ \ }

\newcommand{\seccio}[2]{\par\needspace{16\baselineskip}\bigskip
  \begin{tcolorbox}[colback=navy,colframe=navy,coltext=white,arc=2mm,boxrule=0pt,left=3mm,right=3mm,top=1.2mm,bottom=1.2mm]
  {\large\bfseries #1\quad #2}
  \end{tcolorbox}\vspace{-2pt}}

\newcommand{\exf}[1]{\par\medskip\refstepcounter{exer}\noindent{\color{navy}\bfseries\theexer.}\ \ }
\providecommand{\lt}{<}\providecommand{\gt}{>}
\newcolumntype{C}[1]{>{\centering\arraybackslash}p{#1}}
% quadre de recordatori (només a la fitxa d'exercicis)
\newcommand{\recorda}[1]{\par\smallskip\begin{tcolorbox}[colback=accent!7,colframe=accent!55,arc=2mm,boxrule=0.5pt,left=2.5mm,right=2.5mm,top=1mm,bottom=1mm]\small{\color{accent}\faLightbulb[regular]}\ \textbf{Recorda:}\ #1\end{tcolorbox}}
\newcommand{\res}[1]{#1}\newcommand{\rest}[1]{#1}

\def\fmneed{11}% que el títol d'una secció no quede sol al final del full
% =====================================================================
% Capçalera FentMates de les fitxes: un color per a cada curs.
% Els colors són els de la pàgina inicial (LEVELS a assets/js/materials.js);
% el logo ve de _eines/logo/pdf/ (python3 _eines/logo/genera.py).
%
% Ús, al final del preàmbul de la unitat:
%   % =====================================================================
% Capçalera FentMates de les fitxes: un color per a cada curs.
% Els colors són els de la pàgina inicial (LEVELS a assets/js/materials.js);
% el logo ve de _eines/logo/pdf/ (python3 _eines/logo/genera.py).
%
% Ús, al final del preàmbul de la unitat:
%   % =====================================================================
% Capçalera FentMates de les fitxes: un color per a cada curs.
% Els colors són els de la pàgina inicial (LEVELS a assets/js/materials.js);
% el logo ve de _eines/logo/pdf/ (python3 _eines/logo/genera.py).
%
% Ús, al final del preàmbul de la unitat:
%   \input{../comu/fentmates.tex}
%   \fmcurs{1eso}
% i, al començament del document:
%   \capcalera{Matemàtiques}{1r ESO}{Fitxa d'exercicis · Unitat 2}{Divisibilitat}
% Opcional abans de \input: \def\fmneed{14} (espai mínim abans de cada secció).
% =====================================================================
\makeatletter
\def\fm@logodir{../../logo/pdf}
\providecommand{\fmneed}{10}
\def\fm@cursnom{}

% clau del curs, color i color del text damunt del color
\newcommand{\fm@curs}[3]{\@namedef{fm@color@#1}{#2}\@namedef{fm@text@#1}{#3}}
\fm@curs{1eso}{E6B84A}{navy}
\fm@curs{2eso}{3B82F6}{white}
\fm@curs{3eso}{FF7A50}{white}
\fm@curs{4eso-a}{A78BFA}{white}
\fm@curs{4eso-b}{6D28D9}{white}
\fm@curs{1batx-mi}{9B1C1C}{white}
\fm@curs{1batx-ccss}{EF6B6B}{white}
\fm@curs{2batx-mii}{1E6B47}{white}
\fm@curs{2batx-ccss}{2BB3A3}{white}

\newcommand{\fmcurs}[1]{%
  \@ifundefined{fm@color@#1}{\PackageError{fentmates}{Curs desconegut: #1}{}}{}%
  \edef\fm@tmp{\noexpand\definecolor{curs}{HTML}{\@nameuse{fm@color@#1}}%
               \noexpand\colorlet{curstext}{\@nameuse{fm@text@#1}}}\fm@tmp
  \colorlet{cursfosc}{curs!72!black}%
  \def\fm@logo{\fm@logodir/vertical-#1.pdf}%
  \def\fm@marca{\fm@logodir/marca-#1.pdf}}

% ---------- capçalera de la primera pàgina ----------
% #1 matèria (petit), #2 curs, #3 tipus de document i unitat, #4 títol
\newcommand{\capcalera}[4]{%
  \gdef\fm@cursnom{#2}%
  \par\noindent
  \begin{tcolorbox}[colback=curs,colframe=curs,boxrule=0pt,arc=2.2mm,width=3.5cm,height=2.15cm,
      valign=center,left=3mm,right=2mm,top=1mm,bottom=1mm,nobeforeafter,box align=center]
    \color{curstext}{\footnotesize #1}\par\vspace{1pt}{\fontsize{19}{21}\selectfont\bfseries #2\par}%
  \end{tcolorbox}\hspace{5mm}%
  \begin{minipage}[c]{\dimexpr\linewidth-3.5cm-5mm-2.9cm\relax}
    {\color{cursfosc}\small\bfseries\MakeUppercase{#3}}\par\vspace{2pt}
    {\color{navy}\fontsize{20}{23}\selectfont\bfseries #4\par}
  \end{minipage}\hfill
  \begin{minipage}[c]{2.6cm}\raggedleft\includegraphics[width=2.6cm]{\fm@logo}\end{minipage}\par
  \vspace{2pt}}

% ---------- barres de secció del color del curs ----------
\renewcommand{\seccio}[2]{\par\needspace{\fmneed\baselineskip}\bigskip
  \begin{tcolorbox}[colback=curs,colframe=curs,coltext=curstext,arc=2mm,boxrule=0pt,left=3mm,right=3mm,top=1.2mm,bottom=1.2mm]
  {\large\bfseries #1\quad #2}
  \end{tcolorbox}\vspace{-2pt}}

% ---------- estrelles de dificultat del color del curs (a Batxillerat no se'n posen) ----------
\@ifundefined{estrelles}{}{%
  \renewcommand{\estrelles}[1]{{\color{curs}\ifcase#1\or\faStar\faStar[regular]\faStar[regular]\or\faStar\faStar\faStar[regular]\or\faStar\faStar\faStar\fi}}}

% ---------- capçalera i peu de les pàgines ----------
% primera pàgina: logos institucionals al peu
\fancypagestyle{primera}{\fancyhf{}\renewcommand{\headrulewidth}{0pt}
  \fancyfoot[L]{\raisebox{-0.32\height}{\includegraphics[height=0.78cm]{logo-gva.png}}\hspace{5mm}%
    \raisebox{-0.32\height}{\includegraphics[height=0.78cm]{logo-vermellar.png}}}
  \fancyfoot[R]{\scriptsize\color{muted}\href{https://creativecommons.org/licenses/by-nc-sa/4.0/deed.ca}{\ccbyncsa}\ \ CC BY-NC-SA 4.0 · fentmates.org\qquad\small\thepage}}
% la resta: el símbol a l'esquerra i el curs a la dreta, amb el títol al centre com abans
\fancyhead[L]{\raisebox{-0.3\height}{\includegraphics[height=0.62cm]{\fm@marca}}}
\fancyhead[R]{\small\bfseries\color{cursfosc}\fm@cursnom}
\renewcommand{\headrule}{\color{curs!45}\hrule height \headrulewidth}
\makeatother

%   \fmcurs{1eso}
% i, al començament del document:
%   \capcalera{Matemàtiques}{1r ESO}{Fitxa d'exercicis · Unitat 2}{Divisibilitat}
% Opcional abans de \input: \def\fmneed{14} (espai mínim abans de cada secció).
% =====================================================================
\makeatletter
\def\fm@logodir{../../logo/pdf}
\providecommand{\fmneed}{10}
\def\fm@cursnom{}

% clau del curs, color i color del text damunt del color
\newcommand{\fm@curs}[3]{\@namedef{fm@color@#1}{#2}\@namedef{fm@text@#1}{#3}}
\fm@curs{1eso}{E6B84A}{navy}
\fm@curs{2eso}{3B82F6}{white}
\fm@curs{3eso}{FF7A50}{white}
\fm@curs{4eso-a}{A78BFA}{white}
\fm@curs{4eso-b}{6D28D9}{white}
\fm@curs{1batx-mi}{9B1C1C}{white}
\fm@curs{1batx-ccss}{EF6B6B}{white}
\fm@curs{2batx-mii}{1E6B47}{white}
\fm@curs{2batx-ccss}{2BB3A3}{white}

\newcommand{\fmcurs}[1]{%
  \@ifundefined{fm@color@#1}{\PackageError{fentmates}{Curs desconegut: #1}{}}{}%
  \edef\fm@tmp{\noexpand\definecolor{curs}{HTML}{\@nameuse{fm@color@#1}}%
               \noexpand\colorlet{curstext}{\@nameuse{fm@text@#1}}}\fm@tmp
  \colorlet{cursfosc}{curs!72!black}%
  \def\fm@logo{\fm@logodir/vertical-#1.pdf}%
  \def\fm@marca{\fm@logodir/marca-#1.pdf}}

% ---------- capçalera de la primera pàgina ----------
% #1 matèria (petit), #2 curs, #3 tipus de document i unitat, #4 títol
\newcommand{\capcalera}[4]{%
  \gdef\fm@cursnom{#2}%
  \par\noindent
  \begin{tcolorbox}[colback=curs,colframe=curs,boxrule=0pt,arc=2.2mm,width=3.5cm,height=2.15cm,
      valign=center,left=3mm,right=2mm,top=1mm,bottom=1mm,nobeforeafter,box align=center]
    \color{curstext}{\footnotesize #1}\par\vspace{1pt}{\fontsize{19}{21}\selectfont\bfseries #2\par}%
  \end{tcolorbox}\hspace{5mm}%
  \begin{minipage}[c]{\dimexpr\linewidth-3.5cm-5mm-2.9cm\relax}
    {\color{cursfosc}\small\bfseries\MakeUppercase{#3}}\par\vspace{2pt}
    {\color{navy}\fontsize{20}{23}\selectfont\bfseries #4\par}
  \end{minipage}\hfill
  \begin{minipage}[c]{2.6cm}\raggedleft\includegraphics[width=2.6cm]{\fm@logo}\end{minipage}\par
  \vspace{2pt}}

% ---------- barres de secció del color del curs ----------
\renewcommand{\seccio}[2]{\par\needspace{\fmneed\baselineskip}\bigskip
  \begin{tcolorbox}[colback=curs,colframe=curs,coltext=curstext,arc=2mm,boxrule=0pt,left=3mm,right=3mm,top=1.2mm,bottom=1.2mm]
  {\large\bfseries #1\quad #2}
  \end{tcolorbox}\vspace{-2pt}}

% ---------- estrelles de dificultat del color del curs (a Batxillerat no se'n posen) ----------
\@ifundefined{estrelles}{}{%
  \renewcommand{\estrelles}[1]{{\color{curs}\ifcase#1\or\faStar\faStar[regular]\faStar[regular]\or\faStar\faStar\faStar[regular]\or\faStar\faStar\faStar\fi}}}

% ---------- capçalera i peu de les pàgines ----------
% primera pàgina: logos institucionals al peu
\fancypagestyle{primera}{\fancyhf{}\renewcommand{\headrulewidth}{0pt}
  \fancyfoot[L]{\raisebox{-0.32\height}{\includegraphics[height=0.78cm]{logo-gva.png}}\hspace{5mm}%
    \raisebox{-0.32\height}{\includegraphics[height=0.78cm]{logo-vermellar.png}}}
  \fancyfoot[R]{\scriptsize\color{muted}\href{https://creativecommons.org/licenses/by-nc-sa/4.0/deed.ca}{\ccbyncsa}\ \ CC BY-NC-SA 4.0 · fentmates.org\qquad\small\thepage}}
% la resta: el símbol a l'esquerra i el curs a la dreta, amb el títol al centre com abans
\fancyhead[L]{\raisebox{-0.3\height}{\includegraphics[height=0.62cm]{\fm@marca}}}
\fancyhead[R]{\small\bfseries\color{cursfosc}\fm@cursnom}
\renewcommand{\headrule}{\color{curs!45}\hrule height \headrulewidth}
\makeatother

%   \fmcurs{1eso}
% i, al començament del document:
%   \capcalera{Matemàtiques}{1r ESO}{Fitxa d'exercicis · Unitat 2}{Divisibilitat}
% Opcional abans de \input: \def\fmneed{14} (espai mínim abans de cada secció).
% =====================================================================
\makeatletter
\def\fm@logodir{../../logo/pdf}
\providecommand{\fmneed}{10}
\def\fm@cursnom{}

% clau del curs, color i color del text damunt del color
\newcommand{\fm@curs}[3]{\@namedef{fm@color@#1}{#2}\@namedef{fm@text@#1}{#3}}
\fm@curs{1eso}{E6B84A}{navy}
\fm@curs{2eso}{3B82F6}{white}
\fm@curs{3eso}{FF7A50}{white}
\fm@curs{4eso-a}{A78BFA}{white}
\fm@curs{4eso-b}{6D28D9}{white}
\fm@curs{1batx-mi}{9B1C1C}{white}
\fm@curs{1batx-ccss}{EF6B6B}{white}
\fm@curs{2batx-mii}{1E6B47}{white}
\fm@curs{2batx-ccss}{2BB3A3}{white}

\newcommand{\fmcurs}[1]{%
  \@ifundefined{fm@color@#1}{\PackageError{fentmates}{Curs desconegut: #1}{}}{}%
  \edef\fm@tmp{\noexpand\definecolor{curs}{HTML}{\@nameuse{fm@color@#1}}%
               \noexpand\colorlet{curstext}{\@nameuse{fm@text@#1}}}\fm@tmp
  \colorlet{cursfosc}{curs!72!black}%
  \def\fm@logo{\fm@logodir/vertical-#1.pdf}%
  \def\fm@marca{\fm@logodir/marca-#1.pdf}}

% ---------- capçalera de la primera pàgina ----------
% #1 matèria (petit), #2 curs, #3 tipus de document i unitat, #4 títol
\newcommand{\capcalera}[4]{%
  \gdef\fm@cursnom{#2}%
  \par\noindent
  \begin{tcolorbox}[colback=curs,colframe=curs,boxrule=0pt,arc=2.2mm,width=3.5cm,height=2.15cm,
      valign=center,left=3mm,right=2mm,top=1mm,bottom=1mm,nobeforeafter,box align=center]
    \color{curstext}{\footnotesize #1}\par\vspace{1pt}{\fontsize{19}{21}\selectfont\bfseries #2\par}%
  \end{tcolorbox}\hspace{5mm}%
  \begin{minipage}[c]{\dimexpr\linewidth-3.5cm-5mm-2.9cm\relax}
    {\color{cursfosc}\small\bfseries\MakeUppercase{#3}}\par\vspace{2pt}
    {\color{navy}\fontsize{20}{23}\selectfont\bfseries #4\par}
  \end{minipage}\hfill
  \begin{minipage}[c]{2.6cm}\raggedleft\includegraphics[width=2.6cm]{\fm@logo}\end{minipage}\par
  \vspace{2pt}}

% ---------- barres de secció del color del curs ----------
\renewcommand{\seccio}[2]{\par\needspace{\fmneed\baselineskip}\bigskip
  \begin{tcolorbox}[colback=curs,colframe=curs,coltext=curstext,arc=2mm,boxrule=0pt,left=3mm,right=3mm,top=1.2mm,bottom=1.2mm]
  {\large\bfseries #1\quad #2}
  \end{tcolorbox}\vspace{-2pt}}

% ---------- estrelles de dificultat del color del curs (a Batxillerat no se'n posen) ----------
\@ifundefined{estrelles}{}{%
  \renewcommand{\estrelles}[1]{{\color{curs}\ifcase#1\or\faStar\faStar[regular]\faStar[regular]\or\faStar\faStar\faStar[regular]\or\faStar\faStar\faStar\fi}}}

% ---------- capçalera i peu de les pàgines ----------
% primera pàgina: logos institucionals al peu
\fancypagestyle{primera}{\fancyhf{}\renewcommand{\headrulewidth}{0pt}
  \fancyfoot[L]{\raisebox{-0.32\height}{\includegraphics[height=0.78cm]{logo-gva.png}}\hspace{5mm}%
    \raisebox{-0.32\height}{\includegraphics[height=0.78cm]{logo-vermellar.png}}}
  \fancyfoot[R]{\scriptsize\color{muted}\href{https://creativecommons.org/licenses/by-nc-sa/4.0/deed.ca}{\ccbyncsa}\ \ CC BY-NC-SA 4.0 · fentmates.org\qquad\small\thepage}}
% la resta: el símbol a l'esquerra i el curs a la dreta, amb el títol al centre com abans
\fancyhead[L]{\raisebox{-0.3\height}{\includegraphics[height=0.62cm]{\fm@marca}}}
\fancyhead[R]{\small\bfseries\color{cursfosc}\fm@cursnom}
\renewcommand{\headrule}{\color{curs!45}\hrule height \headrulewidth}
\makeatother

\fmcurs{1batx-mi}

\begin{document}
\thispagestyle{primera}

\capcalera{Matemàtiques I}{1r Batx}{Fitxa d'exercicis · Unitat 3}{\raggedright Sistemes d'equacions i inequacions}
\vspace{10pt}
Nom i cognoms: \rule{8.4cm}{0.4pt}\hfill Grup: \rule{1.6cm}{0.4pt}\hfill Data: \rule{2.2cm}{0.4pt}
\vspace{8pt}

\begin{tcolorbox}[colback=soft,colframe=soft,arc=2mm,left=3mm,right=3mm,top=1.5mm,bottom=1.5mm]
\small Resol els exercicis al quadern mostrant tots els passos. Dona els resultats de manera exacta i, si cal, també aproximats. En els sistemes amb denominadors, radicals o logaritmes, comprova sempre les solucions. Els exercicis 1--28 són de la fitxa del departament; la secció 6 en té de nous per a ampliar i aplicar el que has aprés.
\end{tcolorbox}

% =====================================================
\seccio{1}{Sistemes d'equacions lineals. Mètode de Gauss}

\exf{1} Resol els següents sistemes d'equacions lineals:\begin{tasks}(2)
  \task $\left\{\begin{array}{l}y=\dfrac{x+1}{2}+3\\y=2x+10\end{array}\right.$
  \task $\left\{\begin{array}{l}2(2x+y)-3(3x-2y)=-34\\\dfrac x2-\dfrac y3=2\end{array}\right.$
\end{tasks}

\ex{2} Resol els següents sistemes d'equacions lineals pel mètode de Gauss i classifica'ls:\begin{tasks}(2)
  \task $\left\{\begin{array}{l}x+2y-3z=7\\2x+y-z=6\\3x-y-z=6\end{array}\right.$
  \task $\left\{\begin{array}{l}x+y-z=5\\3x-2y-2z=5\\2x-3y-z=3\end{array}\right.$
  \task $\left\{\begin{array}{l}2x+y+z=5\\2x-y+3z=3\\3x+2y+z=8\end{array}\right.$
  \task $\left\{\begin{array}{l}x+y-2z=0\\2x-3y+3z=4\\5x-5y+4z=8\end{array}\right.$
  \task $\left\{\begin{array}{l}5x+2y-2z=0\\3x-y+3z=0\\8x+y+z=-1\end{array}\right.$
  \task $\left\{\begin{array}{l}x+3y-2z=0\\x-y+3z=0\\2x+2y+z=0\end{array}\right.$
  \task $\left\{\begin{array}{l}2y-z=-1\\5x-y-3z=2\\x-y+2z=3\end{array}\right.$
  \task $\left\{\begin{array}{l}x+2y-2z=2\\3x-3y+z=-14\\5x-y-2z=-15\end{array}\right.$
  \task $\left\{\begin{array}{l}x+y+z=1\\2x-3z=5\\2y+5z=2\end{array}\right.$
  \task $\left\{\begin{array}{l}x+3y-2z+2w=12\\2x-2y-z+w=5\\3x+y-2z-4w=16\\3x-3z-3w=15\end{array}\right.$
\end{tasks}\recorda{Fes zeros davall de la diagonal. L'última fila $(0\ \ 0\ \ m\,|\,n)$ decideix: $m\neq0$ SCD; $m=n=0$ SCI; $m=0,\ n\neq0$ SI.}


% =====================================================
\seccio{2}{Sistemes d'equacions no lineals}

\exf{2} Resol els següents sistemes d'equacions no lineals:\begin{tasks}(3)
  \task $\left\{\begin{array}{l}2x+3y=-4\\x^2-2xy=5\end{array}\right.$
  \task $\left\{\begin{array}{l}\frac1x+\frac2y=4\\10xy=9\end{array}\right.$
  \task $\left\{\begin{array}{l}x^2+2y^2=6\\3x^2-y^2=11\end{array}\right.$
  \task $\left\{\begin{array}{l}2x+y=8\\2x+3y^2=22\end{array}\right.$
  \task $\left\{\begin{array}{l}2x+5y=-8\\xy-3y=4\end{array}\right.$
  \task $\left\{\begin{array}{l}x+y=0\\xy=1\end{array}\right.$
  \task $\left\{\begin{array}{l}2x-y=3\\x^2-y^2=3\end{array}\right.$
  \task $\left\{\begin{array}{l}x+y=11\\\frac1x+\frac6{y+1}=1\end{array}\right.$
  \task $\left\{\begin{array}{l}xy=-3\\x^2+2y^2=19\end{array}\right.$
  \task $\left\{\begin{array}{l}2x^2+3y^2=32\\-3x^2+4y^2=-48\end{array}\right.$
  \task $\left\{\begin{array}{l}\sqrt x+\sqrt y=15\\x-y=105\end{array}\right.$
  \task $\left\{\begin{array}{l}x^2-2xy=16-y^2\\x^2-y^2=72\end{array}\right.$
\end{tasks}

\ex{2} Resol els següents sistemes d'equacions exponencials:\begin{tasks}(2)
  \task $\left\{\begin{array}{l}2^x+5^y=9\\2^{x+2}+5^{y+1}=41\end{array}\right.$
  \task $\left\{\begin{array}{l}2^x+3^y=5\\2^{x+1}-3^y=7\end{array}\right.$
  \task $\left\{\begin{array}{l}2^x-2^y=1016\\2^{x-y}=128\end{array}\right.$
  \task $\left\{\begin{array}{l}3^x+7^y=16\\3^{x-1}-7^{y+2}=-340\end{array}\right.$
  \task $\left\{\begin{array}{l}\dfrac{1}{2^x}+\dfrac{3}{2^y}=\dfrac{x+y}{8}\\x+y=5\end{array}\right.$
\end{tasks}\recorda{Canvi de variable: $u=2^x$, $v=5^y$\ldots\ i resol el sistema lineal en $u$ i $v$.}

\ex{2} Troba les solucions dels següents sistemes d'equacions logarítmiques:\begin{tasks}(2)
  \task $\left\{\begin{array}{l}2\log x-3\log y=7\\\log x+\log y=1\end{array}\right.$
  \task $\left\{\begin{array}{l}x-y=21\\\log x+\log y=2\end{array}\right.$
  \task $\left\{\begin{array}{l}\log x+3\log y=5\\\log x-\log y=3\end{array}\right.$
  \task $\left\{\begin{array}{l}\log_3x=7-\log_3y\\x-2y=27\end{array}\right.$
  \task $\left\{\begin{array}{l}3x+2y=64\\\log x-\log y=1\end{array}\right.$
  \task $\left\{\begin{array}{l}\log x^2-2\log y=6\\2\log x+\log y^2=2\end{array}\right.$
\end{tasks}\recorda{Amb $u=\log x$, $v=\log y$ s'obté un sistema lineal; o bé agrupa els logaritmes. Comprova que $x\gt0$ i $y\gt0$.}


% =====================================================
\seccio{3}{Sistemes d'inequacions amb una incògnita}

\exf{2} Resol els sistemes d'inequacions següents:\begin{tasks}(3)
  \task $\left\{\begin{array}{l}2x-3\le0\\3x+1\gt0\end{array}\right.$
  \task $\left\{\begin{array}{l}x+2\le2\\2x-\frac{x-1}{3}\le x-\frac13\\x^2+2x-8\gt0\end{array}\right.$
  \task $\left\{\begin{array}{l}2x+1\lt x+2\\3x-1\le4x\end{array}\right.$
  \task $\left\{\begin{array}{l}3x-4\gt x\\x\ge2x-1\end{array}\right.$
  \task $\left\{\begin{array}{l}x^2+2x-3\ge0\\2(x-3)\gt-6\end{array}\right.$
  \task $\left\{\begin{array}{l}\dfrac{x-1}{x+2}\gt0\\x^2\ge1\end{array}\right.$
\end{tasks}


% =====================================================
\seccio{4}{Inequacions i sistemes amb dues incògnites}

\exf{1} Resol les següents inequacions amb dues incògnites:\begin{tasks}(2)
  \task $3x+4y\le12$
  \task $-3x+4y\ge4$
  \task $2x+y\lt2$
  \task $x\gt y-2$
\end{tasks}\recorda{Dibuixa la recta (contínua si és $\le,\ge$; discontínua si és $\lt,\gt$) i prova un punt que no hi estiga, com $O(0,0)$.}

\ex{2} Resol els sistemes d'inequacions amb dues incògnites:\begin{tasks}(3)
  \task $\left\{\begin{array}{l}0\le y\lt2\\x+y\le3\\x-y\ge2\\x\ge0,\ y\ge0\end{array}\right.$
  \task $\left\{\begin{array}{l}y\lt-2x+4\\y\ge x\end{array}\right.$
  \task $\left\{\begin{array}{l}y\le2\\x\gt-4\\x-y\lt-3\end{array}\right.$
  \task $\left\{\begin{array}{l}x+y\le6\\2y\le x\\x\ge0\end{array}\right.$
  \task $\left\{\begin{array}{l}2x+y\lt2\\x\gt y-2\\x\gt0,\ y\gt0\end{array}\right.$
  \task $\left\{\begin{array}{l}3\le y\le5\\x-6\le0\\x-4\ge0\end{array}\right.$
\end{tasks}


% =====================================================
\seccio{5}{Problemes de sistemes i inequacions}

\exf{1} En una granja hi ha 250 animals entre conills i gallines. Si el nombre total de potes és 800, quants animals hi ha de cada tipus?

\ex{3} Hem necessitat 10 dm$^2$ de cartó per a construir una caixa amb tapa, de base quadrada i 2 dm$^3$ de volum. Quines són les dimensions de la caixa si el costat de la base és un nombre enter de decímetres?

\ex{2} En una caixa registradora hi ha bitllets de 50\,€, 100\,€ i 200\,€: 21 bitllets en total i 1\,800\,€. El nombre de bitllets de 50\,€ és el quíntuple dels de 200\,€. Quants n'hi ha de cada classe?

\ex{2} Tres néts volen fer un regal de 60\,€ a la seua àvia: Lluís, el major, aporta el triple del que aporten els altres dos junts, i Carme aporta 3\,€ per cada 2 que aporta Pere. Quant aporta cadascun?

\ex{2} Calcula quantes bicicletes, tricicles i cotxes té una botiga de joguets si en total hi ha 50 vehicles, el nombre de bicicletes i tricicles quadruplica el de cotxes i hi ha 140 rodes.

\ex{2} Per a un espectacle hi ha entrades A, B i C, de 10, 15 i 30\,€. La recaptació va ser de 775\,€, es van vendre 5 entrades més del tipus B que del C i el 20\,\% de les entrades eren del tipus C. Quantes se'n van vendre de cada tipus?

\ex{2} Un caixer només té bitllets de 10, 20 i 50\,€. Traiem 290\,€ i ens dona exactament 8 bitllets; el nombre de bitllets de 10\,€ és el doble dels de 20\,€. Planteja i resol el sistema.

\ex{2} Un caixer automàtic conté 800 bitllets de 10, 20 i 50\,€ amb un import de 21\,000\,€. El nombre de bitllets de 10\,€ és igual que el de 20\,€ i 50\,€ junts. Quants bitllets hi ha de cada tipus?

\ex{3} En un cinema s'han venut 1405 entrades i s'han recaptat 7920\,€. L'entrada normal costa 6\,€, la del dia de l'espectador 4\,€ i la dels jubilats 3\,€. La recaptació de les entrades de preu reduït és el 10\,\% de la de les entrades normals. Quantes entrades de cada tipus s'han venut?

\ex{2} Un comerciant ven rellotges A a 200\,€, B a 500\,€ i C a 250\,€. Un mes en ven 200; els del tipus B són tants com els de A i C junts, i la recaptació és de 73\,500\,€. Quants n'ha venut de cada tipus?

\ex{2} Una persona compra memòries externes, lectors de llibres electrònics i tauletes gràfiques a 100, 120 i 150\,€ la unitat. Paga 1160\,€ per 9 unitats i compra una tauleta més que lectors. Quantes unitats de cada producte compra?

\ex{2} Després d'un descompte del 10\,\% en cada preu, s'ha pagat 3,96\,€ per un retolador, un quadern i una carpeta. El quadern val la meitat que el retolador i la carpeta, el que val el quadern més el 20\,\% del retolador. Calcula el preu original de cada objecte.

\ex{3} Una persona inverteix 7000\,€ en accions de les empreses A i B i en un dipòsit a l'1\,\%. Amb les accions obté un 5\,\% (A) i un 3\,\% (B), i el benefici total és de 202\,€. Els diners en accions superen en 2600\,€ els del dipòsit. Quant va destinar a cada inversió?

\ex{2} Una agència de viatges ven 60 bitllets a Londres, París i Roma. Els de París són el doble dels altres dos destins junts i a Roma se n'emeten dos més que la meitat dels de Londres. Quants bitllets s'han venut per a cada destí?

\ex{3} Un comerciant compra 200 kg de bresquilles, 100 de pomes i 300 de peres. Els ven apujant un 25\,\% el preu de les bresquilles i les pomes i un 40\,\% el de les peres, i obté 1087\,€, dels quals 257 són de benefici. Si el quilo de bresquilles li va costar 50 cèntims més que el de peres, quin va ser el preu de compra de cada fruita?

\ex{3} Un restaurant ofereix desdejunis a 4\,€, dinars a 8\,€ i sopars a 10\,€. L'últim dissabte es van servir tants dinars com desdejunis i sopars junts, la recaptació total va ser de 1116\,€ i la dels dinars va superar la dels sopars en 156\,€.\begin{tasks}(1)
  \task Quants desdejunis, dinars i sopars es van servir?
  \task Quin benefici es va obtindre si els guanys són de 2,5\,€ per desdejuni, 4\,€ per dinar i 5\,€ per sopar?
\end{tasks}

\ex{2} Un estudiant va traure 7,5 punts en un examen de tres preguntes. En la tercera va traure un punt més que en la segona, i els punts de la primera quintuplicaven la diferència entre els de la tercera i la primera. Quina puntuació va traure en cada pregunta?

\ex{3} Un inversor reparteix 42\,000\,€ entre un compte d'estalvis al 5\,\%, un dipòsit a termini fix al 7\,\% i uns bons al 9\,\% anual. En un any obté 2600\,€ d'interessos, i els del compte d'estalvis són 200\,€ menys que la suma dels de les altres dues inversions. Quant va invertir en cada producte?

\ex{2} La llargària d'un rectangle és 1 cm més que el doble de l'amplària. Volem que el perímetre siga com a mínim de 20 cm i com a màxim de 30 cm. Entre quins valors pot estar l'amplària?

\ex{2} Un artesà fabrica collarets i polseres. Un collaret necessita 2 hores de treball i 1 unitat de material, i una polsera, 1 hora i 2 unitats de material. Cada setmana disposa, com a màxim, de 40 hores i 50 unitats de material.\begin{tasks}(1)
  \task Planteja el sistema d'inequacions i representa la regió de les possibles combinacions de collarets i polseres.
  \task Pot fabricar 15 collarets i 12 polseres en una setmana? I 10 collarets i 20 polseres?
\end{tasks}


% =====================================================
\seccio{6}{Ampliació i problemes}

\exf{3} Classifica el sistema segons els valors de $a$ i resol-lo quan siga compatible: \ $\left\{\begin{array}{l}x+y+z=3\\x-y+z=1\\2x+az=4\end{array}\right.$

\ex{2} Un rectangle té 10 cm de perímetre i 6 cm$^2$ d'àrea. Quines en són les dimensions?

\ex{2} Fa 2 anys l'edat d'un pare era el triple de la del fill i d'ací a 10 anys serà el doble. Quants anys tenen ara?

\ex{3} Escriu un sistema d'inequacions la solució del qual siga el triangle de vèrtexs $(0,0)$, $(4,0)$ i $(0,2)$, amb les vores incloses.

\vfill
\begin{center}\color{muted}\small\faCheckCircle[regular]\ \ Ànim! Comprova els teus resultats amb les eines interactives dels apunts de la unitat.\end{center}

\end{document}
